\documentclass[Markos_Delaportas_1316.tex]{subfiles}

\begin{document}

\section{Conclusion and Future Work}

\subsection{Summary of Thesis Findings}
This thesis has systematically investigated the feasibility,
theoretical scaling, and physical
limitations of Over-the-Air (OTA) computation. By developing a
comprehensive cross-platform
modeling methodology---spanning mathematical derivation,
discrete-time MATLAB simulation,
continuous-time Simulink architectures, and Software-Defined Radio
(SDR) realization---the
research bridged the significant gap between idealized analog
computation theory and practical
deployment.

The fundamental analysis confirmed that under ideal conditions,
\gls{ota} computation achieves
unprecedented spectral efficiency by exploiting the superposition
property of the wireless
Multiple-Access Channel. However, the rigorous evaluation of channel
impairments demonstrated
that the system's performance is severely bottlenecked by
uncoordinated transmitter dynamics.
Specifically, multipath propagation and fractional timing delays
induce Intersymbol Interference
(ISI) that manifests as an irreducible error floor. Furthermore,
asynchronous Carrier Frequency
Offsets (CFO) cause catastrophic coherent demodulation failure due to
beat-frequency oscillations.
While non-coherent envelope detection provides a robust fallback, it
suffers from vector magnitude
fading. The findings conclusively show that traditional
point-to-point equalization schemes (e.g.,
standard Costas loops) are insufficient for multi-user aggregated
signals, necessitating paradigm
shifts in how synchronization and filtering are handled for \gls{ota} networks.

\subsection{Future Research Directions}
The insights gained from the physical degradation models open several
critical pathways for
future research in \gls{ota} computation and federated edge learning
\cite{xiao_heterogeneous_2025, zhu_over--air_2024}:
\begin{itemize}
    \item \textbf{Adaptive and Non-Linear Equalization:}
        Transitioning from static recovery loops to dynamic, adaptive
        algorithms is paramount. Future work will explore the
        deployment of Unscented Kalman Filters (UKF) and Extended
        Kalman Filters (EKF) for the real-time tracking and
        compensation of non-linear impairments directly from the
        aggregated waveform.
    \item \textbf{Multi-Antenna and MIMO Systems:} Expanding the
        single-antenna architecture to Multiple-Input Multiple-Output
        (MIMO) systems offers spatial degrees of freedom to combat
        fading. Implementing spatial equalization and receive
        beamforming \cite{sahin_survey_2023} could significantly
        mitigate the destructive interference caused by asynchronous
        phase offsets.
    \item \textbf{Large-Scale Federated Learning Deployments:} As
        \gls{ota} computation proves to be a natural fit for the
        aggregation step in Federated Learning, future physical
        deployments should evaluate the impact of the quantified
        \gls{ota} error floors on the convergence rates of
        distributed neural network training.
\end{itemize}

\ifSubfilesClassLoaded{
    \clearpage
    \printbibliography
}{}

\end{document}
